\section{Canonical-content inversion and the retained bilinear form}
\label{sec:tpc32-canonical-factorization}

We now place the canonical determinant phase inside the exact cutoff
shell.  The central point is elementary but useful: exact target
content can be imposed by a common-divisor M\"obius projector, and the
normalized determinant character then splits into one character on
each row.  The fixed residue factor, physical weight, and pair mask are
not discarded in this operation.

\subsection{Primitive common divisors}

For two physical rows \(m=m_\alpha\) and \(n=m_\gamma\), abbreviate
\begin{equation}
 N_m=mj+h,\qquad N_n=nj+h,
 \qquad G_{m,n}(j)=(N_m,N_n).
 \label{eq:tpc32-two-targets-content}
\end{equation}
When \(m\ne n\), the canonical normalized determinant of TPC-31 is
\begin{equation}
 \Delta^\#_{m,n}(j)=\frac{m-n}{G_{m,n}(j)}.
 \label{eq:tpc32-canonical-determinant}
\end{equation}

\begin{lemma}[Every common target divisor divides the row difference]
\label{lem:tpc32-common-divisor-row-difference}
Under \eqref{eq:tpc32-primitive-target-support}, if
\begin{equation}
 t\mid N_m,\qquad t\mid N_n,
 \label{eq:tpc32-common-target-divisor}
\end{equation}
then
\begin{equation}
 (t,mnhj)=1,\qquad t\mid m-n.
 \label{eq:tpc32-common-divisor-properties}
\end{equation}
In particular
\begin{equation}
 G_{m,n}(j)\mid m-n,
 \label{eq:tpc32-content-divides-difference}
\end{equation}
so \eqref{eq:tpc32-canonical-determinant} is a nonzero integer.
\end{lemma}

\begin{proof}
Target primitivity \eqref{eq:tpc32-target-divisor-primitivity} applied
on each row gives \((t,mnH)=1\).  If a prime \(p\) divided both \(t\)
and \(j\), then \(p\mid N_m-mj=h\), hence \(p\mid H\), a
contradiction.  Thus \((t,j)=1\).  Finally
\[
 t\mid N_m-N_n=(m-n)j,
\]
and \(j\) may be cancelled modulo \(t\).  Taking \(t=G_{m,n}(j)\)
proves the last assertion.  This is the primitive target-gcd geometry
used in TPC-30 and TPC-31
\citep{WangTPC30,WangTPC31}.
\end{proof}

\subsection{Exact-content M\"obius projectors}

The following formulation is convenient for both cumulative and
dyadic content ranges.

\begin{definition}[Content projector]
\label{def:tpc32-content-projector}
For a finite set \(\mathscr C\subset\N\), define
\begin{equation}
 \eta_{\mathscr C}(t)
 :=\sum_{\substack{c\mid t\\c\in\mathscr C}}
 \mu(t/c).
 \label{eq:tpc32-content-projector-coefficient}
\end{equation}
\end{definition}

\begin{proposition}[Exact gcd M\"obius inversion]
\label{prop:tpc32-gcd-mobius-inversion}
For all positive integers \(N_1,N_2\), with \(G=(N_1,N_2)\),
\begin{equation}
 \boxed{
 \one_{G\in\mathscr C}
 =\sum_{\substack{t\mid N_1\\t\mid N_2}}
   \eta_{\mathscr C}(t).}
 \label{eq:tpc32-gcd-mobius-inversion}
\end{equation}
Equivalently,
\begin{equation}
 \one_{G=c}
 =\sum_{k\ge1}\mu(k)
   \one_{ck\mid N_1}\one_{ck\mid N_2}
 \qquad(c\ge1).
 \label{eq:tpc32-exact-content-inversion}
\end{equation}
All sums in these identities are finite.
\end{proposition}

\begin{proof}
Interchanging the divisor sums gives
\begin{align*}
 \sum_{t\mid G}\eta_{\mathscr C}(t)
 &=\sum_{\substack{c\mid G\\c\in\mathscr C}}
   \sum_{k\mid G/c}\mu(k).
\end{align*}
The inner sum is one exactly when \(G/c=1\), and is zero otherwise.
This proves \eqref{eq:tpc32-gcd-mobius-inversion}.  Taking
\(\mathscr C=\{c\}\) gives
\eqref{eq:tpc32-exact-content-inversion}.
\end{proof}

\begin{remark}[Size and sign]
The coefficient \(\eta_{\mathscr C}\) is signed and need not be
multiplicative.  The general pointwise estimate
\begin{equation}
 |\eta_{\mathscr C}(t)|\le\tau(t)
 \label{eq:tpc32-projector-trivial-bound}
\end{equation}
does not by itself give cancellation after summation over \(t\).
The value of the projector is its exact reassembly, not an individual
smallness property of its coefficients.
\end{remark}

\subsection{Canonical determinant phases}

Write
\begin{equation}
 e_q(x)=\exp(2\pi i x/q),\qquad q\ge1.
 \label{eq:tpc32-additive-character}
\end{equation}

\begin{theorem}[Exact canonical phase splitting]
\label{thm:tpc32-canonical-phase-splitting}
Let \(q\ge1\), \(r\in\Z\), let \(m\ne n\), and let
\(\mathscr C\subset\N\) be finite.  Under the hypotheses of
\cref{lem:tpc32-common-divisor-row-difference},
\begin{align}
 &\one_{G_{m,n}(j)\in\mathscr C}
 e_q\!\left(-r\Delta^\#_{m,n}(j)\right)
 \notag\\
 &\quad={}
 \sum_{c\in\mathscr C}\sum_{k\ge1}\mu(k)
 \one_{ck\mid N_m}\one_{ck\mid N_n}
 e_{cq}\!\left(-r(m-n)\right).
 \label{eq:tpc32-canonical-phase-inversion}
\end{align}
For every summand on the right,
\begin{equation}
 \boxed{
 e_{cq}\!\left(-r(m-n)\right)
 =e_{cq}(-rm)e_{cq}(rn).}
 \label{eq:tpc32-canonical-phase-split}
\end{equation}
Neither identity requires \((c,q)=1\).
\end{theorem}

\begin{proof}
Fix \(c\in\mathscr C\).  If \(ck\) divides both targets, then
\cref{lem:tpc32-common-divisor-row-difference} gives \(c\mid m-n\),
so every displayed character is well defined.  For the actual content
\(G=G_{m,n}(j)\), the contribution of this \(c\) is zero unless
\(c\mid G\); when \(c\mid G\), it equals
\[
 e_{cq}\!\left(-r(m-n)\right)
 \sum_{k\mid G/c}\mu(k).
\]
M\"obius inversion kills it unless \(G=c\).  In the surviving case,
\[
 e_{cq}\!\left(-r(m-n)\right)
 =e_q\!\left(-r\frac{m-n}{G}\right),
\]
which proves \eqref{eq:tpc32-canonical-phase-inversion}.  Equation
\eqref{eq:tpc32-canonical-phase-split} is the additive-character
identity obtained by expanding \(m-n\).  No modular inverse of \(c\)
has been used, so no coprimality between \(c\) and \(q\) is present.
\end{proof}

\begin{remark}[Why the denominator is \(cq\)]
It would be incorrect to replace division by \(c\) with multiplication
by an inverse of \(c\pmod q\).  The phase is lifted from modulus \(q\)
to modulus \(cq\); this is precisely what makes
\eqref{eq:tpc32-canonical-phase-split} valid even when \((c,q)>1\).
\end{remark}

\subsection{The complete retained-multiplier factorization}

For a finite content family \(\mathscr C\), define the matched
determinant transform
\begin{align}
 \widehat{\cS}^{\mathrm{sh}}_{\mathscr C,q}(r)
 :={}&\sum_{\substack{\alpha,\gamma\\m_\alpha\ne m_\gamma}}
 \gamma_\alpha^{(1)}\gamma_\gamma^{(2)}
 \sum_j\mathfrak A_{\alpha,\gamma}(j)
 \cK^{\mathrm{sh}}_{\alpha,\gamma}(j)
 \notag\\
 &\qquad\times
 \one_{G_{m_\alpha,m_\gamma}(j)\in\mathscr C}
 e_q\!\left(-r\Delta^\#_{m_\alpha,m_\gamma}(j)\right).
 \label{eq:tpc32-matched-determinant-transform}
\end{align}
This is the sum of the transforms of all three raw channels, before
any triangle inequality.

Put
\begin{equation}
 \epsilon_{U_0}=1,\qquad \epsilon_T=-1.
 \label{eq:tpc32-cutoff-signs}
\end{equation}
For \(Y\in\{U_0,T\}\), define the two one-row factors
\begin{align}
 X^{(1)}_{\alpha;c,k,Y,r}(j)
 :={}&\gamma_\alpha^{(1)}
 \one_{ck\mid N_\alpha(j)}
 \mathsf A_{m_\alpha,Y}(j)
 e_{cq}(-rm_\alpha),
 \label{eq:tpc32-left-row-factor}\\
 X^{(2)}_{\gamma;c,k,Y,r}(j)
 :={}&\gamma_\gamma^{(2)}
 \one_{ck\mid N_\gamma(j)}
 \mathsf A_{m_\gamma,Y}(j)
 e_{cq}(rm_\gamma).
 \label{eq:tpc32-right-row-factor}
\end{align}

\begin{theorem}[Matched bilinear factorization with joint multiplier]
\label{thm:tpc32-retained-bilinear-factorization}
Under the physical hypotheses of
\cref{sec:tpc32-physical-shell}, for every finite
\(\mathscr C\subset\N\), every \(q\ge1\), and every \(r\in\Z\),
\begin{align}
 \widehat{\cS}^{\mathrm{sh}}_{\mathscr C,q}(r)
 ={}&\sum_{c\in\mathscr C}\sum_{k\ge1}\mu(k)
 \sum_j\sum_{Y\in\{U_0,T\}}\epsilon_Y
 \notag\\
 &\times\sum_{\substack{\alpha,\gamma\\m_\alpha\ne m_\gamma}}
 \mathfrak A_{\alpha,\gamma}(j)
 X^{(1)}_{\alpha;c,k,Y,r}(j)
 X^{(2)}_{\gamma;c,k,Y,r}(j).
 \label{eq:tpc32-retained-bilinear-factorization}
\end{align}
The sums over \(k\) are finite.  Formula
\eqref{eq:tpc32-retained-bilinear-factorization} retains the exact
rank--two cutoff difference and the complete joint physical
multiplier.
\end{theorem}

\begin{proof}
Insert the cutoff difference
\eqref{eq:tpc32-exact-cutoff-shell} into
\eqref{eq:tpc32-matched-determinant-transform}.  Apply
\cref{thm:tpc32-canonical-phase-splitting} to the content indicator and
phase.  The two row characters in
\eqref{eq:tpc32-canonical-phase-split}, the two common-divisibility
indicators, the two row coefficients, and the two prefixes are exactly
the factors in
\eqref{eq:tpc32-left-row-factor}--
\eqref{eq:tpc32-right-row-factor}.  The joint multiplier is never
altered.  Since \(ck\) divides two targets of size
\(O_{\mathscr D}(X)\), one has \(ck\ll_{\mathscr D}X\), proving
finiteness.
\end{proof}

For fixed \(c,k,Y,r,j\), let
\(M_j=(\mathfrak A_{\alpha,\gamma}(j))_{\alpha,\gamma}\).  The innermost
sum in \eqref{eq:tpc32-retained-bilinear-factorization} is the bilinear
form
\begin{equation}
 (X^{(1)}_{c,k,Y,r}(j))^{\mathsf T}
 M_jX^{(2)}_{c,k,Y,r}(j).
 \label{eq:tpc32-matrix-bilinear-form}
\end{equation}
No complex conjugation is implicit.  This notation makes the remaining
boundary explicit: entrywise boundedness of \(M_j\) does not imply a
bounded operator norm or a low-cost projective decomposition.

\begin{corollary}[Conditional full product representation]
\label{cor:tpc32-conditional-projective-factorization}
Suppose, on a particular physical packet, that the retained multiplier
has an explicit representation
\begin{equation}
 \mathfrak A_{\alpha,\gamma}(j)
 =\int_\Omega
 p_{\omega,j}(\alpha)q_{\omega,j}(\gamma)
 \,d\nu(\omega),
 \label{eq:tpc32-joint-projective-representation}
\end{equation}
with absolute convergence sufficient to interchange all finite sums.
Then the row-pair sum in
\eqref{eq:tpc32-retained-bilinear-factorization} equals
\begin{align}
 \int_\Omega
 &\left(\sum_\alpha
 p_{\omega,j}(\alpha)
 X^{(1)}_{\alpha;c,k,Y,r}(j)\right)
 \notag\\
 &\times
 \left(\sum_\gamma
 q_{\omega,j}(\gamma)
 X^{(2)}_{\gamma;c,k,Y,r}(j)\right)
 \,d\nu(\omega).
 \label{eq:tpc32-full-product-representation}
\end{align}
\end{corollary}

\begin{proof}
Substitute \eqref{eq:tpc32-joint-projective-representation} and
interchange the absolutely convergent integral and sums.
\end{proof}

The fixed-residue and physical smooth factors have such bounded-cost
representations in the inherited packet \citep{WangTPC25}.  The
generic pair mask has not been proved to satisfy one with a sufficiently
small projective norm.  Thus
\cref{cor:tpc32-conditional-projective-factorization} is an exact
conditional interface, not an estimate asserted from
\eqref{eq:tpc32-joint-multiplier-bound}.

\subsection{Dyadic projectors and exact zero-frequency reassembly}

For \(C\ge1\), define the cumulative and dyadic coefficients
\begin{align}
 \eta_{\le C}(t)
 &:=\sum_{\substack{c\mid t\\c\le C}}\mu(t/c),
 \label{eq:tpc32-cumulative-projector}\\
 \eta_{(C,2C]}(t)
 &:=\sum_{\substack{c\mid t\\C<c\le2C}}\mu(t/c).
 \label{eq:tpc32-dyadic-projector}
\end{align}
To match the cumulative notation used elsewhere in the paper, we also
write
\begin{equation}
 \lambda_C(t):=\eta_{\le C}(t).
 \label{eq:tpc32-lambda-projector-alias}
\end{equation}

\begin{lemma}[The first dyadic band]
\label{lem:tpc32-first-dyadic-projector-band}
For every \(C\ge1\),
\begin{equation}
 \boxed{
 \lambda_C(1)=1,\qquad
 \lambda_C(t)=0\ (2\le t\le C),\qquad
 \lambda_C(t)=-1\ (C<t\le2C).}
 \label{eq:tpc32-first-dyadic-projector-band}
\end{equation}
\end{lemma}

\begin{proof}
The first equality is immediate.  If \(2\le t\le C\), every divisor of
\(t\) is included in the definition, so
\[
 \lambda_C(t)=\sum_{c\mid t}\mu(t/c)
 =\sum_{d\mid t}\mu(d)=0.
\]
If \(C<t\le2C\), every proper divisor of \(t\) is at most \(t/2\le C\),
whereas \(c=t\) is omitted.  Therefore
\[
 \lambda_C(t)=\sum_{\substack{c\mid t\\c<t}}\mu(t/c)
 =\sum_{\substack{d\mid t\\d>1}}\mu(d)
 =-\mu(1)=-1.
\]
\end{proof}

\begin{corollary}[Exact cumulative and dyadic content projectors]
\label{cor:tpc32-dyadic-content-projector}
For \(G=(N_1,N_2)\),
\begin{align}
 \one_{G\le C}
 &=\sum_{t\mid N_1,\,t\mid N_2}\eta_{\le C}(t),
 \label{eq:tpc32-cumulative-projector-identity}\\
 \one_{C<G\le2C}
 &=\sum_{t\mid N_1,\,t\mid N_2}\eta_{(C,2C]}(t).
 \label{eq:tpc32-dyadic-projector-identity}
\end{align}
Moreover,
\begin{equation}
 \eta_{(C,2C]}(t)
 =\eta_{\le2C}(t)-\eta_{\le C}(t).
 \label{eq:tpc32-dyadic-projector-difference}
\end{equation}
Hence cumulative projectors telescope exactly into dyadic content
layers, with an optionally truncated final layer.
\end{corollary}

\begin{proof}
Apply \cref{prop:tpc32-gcd-mobius-inversion} to
\(\mathscr C=[1,C]\cap\N\) and
\(\mathscr C=(C,2C]\cap\N\).  The coefficient identity follows
directly from the definitions.
\end{proof}

At determinant frequency zero, the exact-content labels may be
reassembled into one common-divisor coefficient.  Indeed, setting
\(r=0\) in \eqref{eq:tpc32-canonical-phase-inversion} and writing
\(t=ck\) gives
\begin{equation}
 \one_{G\le C}
 =\sum_{t\mid N_1,\,t\mid N_2}\eta_{\le C}(t).
\end{equation}
The term \(t=1\) is distinguished because
\begin{equation}
 \eta_{\le C}(1)=1.
 \label{eq:tpc32-t-one-coefficient}
\end{equation}

\begin{proposition}[\(t=1\) and large-content reassembly]
\label{prop:tpc32-t-one-large-reassembly}
For every positive integer \(G\) and every \(C\ge1\),
\begin{equation}
 \boxed{
 \sum_{\substack{t\mid G\\t>1}}\eta_{\le C}(t)
 =-\one_{G>C}.}
 \label{eq:tpc32-nontrivial-projector-reassembly}
\end{equation}
Thus the \(t=1\) term of the cumulative projector is the unrestricted
amplitude, while all \(t>1\) terms reassemble exactly to the negative
large-content amplitude.
\end{proposition}

\begin{proof}
By \eqref{eq:tpc32-cumulative-projector-identity},
\(\sum_{t\mid G}\eta_{\le C}(t)=\one_{G\le C}\).  Subtract the
\(t=1\) value \eqref{eq:tpc32-t-one-coefficient} and use
\(\one_{G\le C}-1=-\one_{G>C}\).
\end{proof}

To state the consequence for the physical shell, let
\begin{align}
 \cS_{\mathrm{full}}^{\mathrm{sh}}
 &:={}
 \cS[\cK^{\mathrm{sh}};\textup{true}],
 \label{eq:tpc32-full-raw-shell}\\
 \cS_{>C}^{\mathrm{sh}}
 &:={}
 \cS[\cK^{\mathrm{sh}};G_{m_\alpha,m_\gamma}(j)>C].
 \label{eq:tpc32-large-content-raw-shell}
\end{align}

\begin{corollary}[Exact auxiliary-zero reassembly]
\label{cor:tpc32-zero-reassembly}
For every auxiliary modulus \(q\ge1\),
\begin{equation}
 \boxed{
 \widehat{\cS}^{\mathrm{sh}}_{[1,C],q}(0)
 =\cS_{\mathrm{full}}^{\mathrm{sh}}
  -\cS_{>C}^{\mathrm{sh}}.}
 \label{eq:tpc32-zero-reassembly}
\end{equation}
More precisely, the \(t=1\) term in the common-divisor expansion is
\(\cS_{\mathrm{full}}^{\mathrm{sh}}\), and the sum of all \(t>1\)
terms is \(-\cS_{>C}^{\mathrm{sh}}\).
\end{corollary}

\begin{proof}
At \(r=0\), the phase in
\eqref{eq:tpc32-matched-determinant-transform} is one.  Apply
\cref{prop:tpc32-t-one-large-reassembly} pointwise to each physical
row--orbit triple, then sum with the original signed coefficient and
the retained multiplier.
\end{proof}

TPC-30 proves, uniformly for the sum of all three raw channels,
\begin{equation}
 |\cS_{>C}^{\mathrm{sh}}|
 \ll_{\eps,h,\mathscr D}
 X^\eps\left(Q^2+\frac{XQ}{C}\right)
 \ll_{\mathscr D}X^\eps XQ
 \left(\frac1J+\frac1C\right).
 \label{eq:tpc32-inherited-large-content-bound}
\end{equation}
Combining this bound with
\eqref{eq:tpc32-zero-reassembly} gives
\begin{equation}
 \left|
 \widehat{\cS}^{\mathrm{sh}}_{[1,C],q}(0)
 -\cS_{\mathrm{full}}^{\mathrm{sh}}
 \right|
 \ll X^\eps XQ\left(\frac1J+\frac1C\right).
 \label{eq:tpc32-zero-full-equivalence}
\end{equation}
This is an equivalence up to an already controlled large-content
error, not an estimate for the unrestricted shell.

For the complete calibrated cutoff difference, the drift estimate
\eqref{eq:tpc32-drift-boundary} yields the corresponding boundary
\begin{align}
 \cK^{\mathrm{cal}}_{\alpha,\gamma}(j)
 &:=
 P_{m_\alpha,U_0}(j)P_{m_\gamma,U_0}(j)
 -P_{m_\alpha,T}(j)P_{m_\gamma,T}(j),
 \label{eq:tpc32-calibrated-shell-kernel}\\
 \cS_{\mathrm{full}}^{\mathrm{cal}}
 &:=
 \cS[\cK^{\mathrm{cal}};\textup{true}],
 \notag\\
 \widehat{\cS}^{\mathrm{cal}}_{[1,C],q}(0)
 &:=
 \cS[\cK^{\mathrm{cal}};G_{m_\alpha,m_\gamma}(j)\le C].
 \label{eq:tpc32-calibrated-zero-definition}
\end{align}
The notation retains \(q\) to emphasize that this is the zero
coefficient of the same auxiliary residue-label transform; its value
is independent of \(q\).  One then has
\begin{equation}
 \boxed{
 \left|
 \widehat{\cS}^{\mathrm{cal}}_{[1,C],q}(0)
 -\cS_{\mathrm{full}}^{\mathrm{cal}}
 \right|
 \ll_{\eps,h,\mathscr D}
 X^\eps XQ
 \left(\frac1J+\frac1C+\frac1L\right).}
 \label{eq:tpc32-calibrated-zero-full-equivalence}
\end{equation}
Here both calibrated quantities are defined using the two-prefix
difference in \eqref{eq:tpc32-calibrated-cutoff}; the \(L^{-1}\) term
accounts for the two drift legs uniformly under either the small- or
large-content restriction.

\begin{remark}[The zero-frequency gate is not removed]
Equation \eqref{eq:tpc32-zero-full-equivalence} shows that controlling
the auxiliary residue-label zero coefficient is, within the inherited
large-content margin, the same problem as controlling the unrestricted
matched shell.  The exact M\"obius projector exposes this fact but does
not make the \(t=1\) contribution oscillatory.  Bounds for nonzero
determinant frequencies, or termwise absolute bounds for
\(\eta_{\le C}(t)\), do not imply the missing estimate.
\end{remark}
